\section{Introduction}

The deployment of autonomous multi-agent systems into consequential environments has exposed a structural failure mode that policy alone cannot address. As agent populations grow recursively—spawning children, modifying state, executing decisions with real-world consequences—the question of who authorized each action, through what chain of delegation, and with what scope constraints, has become the central unsolved problem in agentic infrastructure. The failure mode is called \emph{autonomous drift}: the condition in which an agent continues to operate with no runtime-verifiable authorization path to a human principal. The agent functions normally, continues to spawn, continues to make decisions, yet the delegation chain that originally authorized its provisioning has been severed, obscured, or was never structurally established. The agent persists in operational state, accumulating consequences that no human can attribute or reverse. This is not a conventional system malfunction—the system continues to operate. The failure is structural: accountability has been severed at the architectural level, not at the policy level.

Autonomous drift arises from a specific combination of three conditions that conventional architectures accept as benign but which, taken together, make drift not merely possible but structurally inevitable. First, the parent pointer that links a child agent to its parent is held in mutable state, allowing the child to rewrite its own ancestry either explicitly or through a cascading state corruption. Second, the operating scope of each agent is not enforced as a descendant refinement of its parent's scope, permitting each generation to expand its authority beyond what its parent was authorized to delegate. Third, there is no structural depth bound enforced by an independent execution context—depth limits, when they exist at all, are implemented as runtime parameters that the agent itself can inspect and circumvent. An agent near completion of a high-value task has a structural incentive to suppress escalation, exceed its token budget, or absorb exceptions rather than surface them. Policy-based constraints cannot close this gap because they operate through the same execution context that benefits from their violation.

The standard responses to this problem—time-to-live (TTL) expiration and depth limiting—address symptoms rather than the structural root cause. TTL assumes that an agent that has run long enough has exhausted its legitimate authorization. This is a statistical argument applied to an architectural problem: a drift-compliant agent does not need to run longer than the TTL to cause harm; it only needs to exceed its provisioned scope within the TTL window. Moreover, TTL is enforced by the same runtime that the agent inhabits, making it bypassable through clock manipulation, process forking, or coordinated TTL suppression across a spawn chain. Depth limiting is similarly insufficient: it caps the number of spawn generations but does not constrain the scope or authority of any individual agent, leaving the drift problem intact within any depth bound. More fundamentally, neither TTL nor depth limits address the mutable parent pointer, which is the structural enabler of drift. If a child agent can rewrite the reference to its parent after being spawned, no depth bound or TTL can prevent it from presenting a false ancestry and continuing to operate indefinitely under a fabricated authorization chain.

The core insight of the Recursive Spawning Protocol is that autonomous drift is structurally impossible when the parent pointer $\alpha(n)$ is immutable—established once at provisioning and unalterable by any actor under any circumstances. An immutable parent pointer converts the delegation chain from a mutable runtime assumption into a cryptographically enforceable structural artifact. Every child carries its parent reference as a first-class property of its own node identity, not as a pointer in mutable memory. The chain is therefore traceable at runtime by an independent enforcement layer, not reconstructable only after the fact from logs that the agent itself controls. This is not a stronger policy—it is a different kind of mechanism entirely. Policy operates within the agent's execution context. Structure operates outside it.

We define drift formally as follows. Let $\mathcal{N}(t)$ be the set of active nodes at time $t$, and for each node $n \in \mathcal{N}(t)$, let $\alpha(n)$ denote its parent node (or $\text{nil}$ if $n$ is a root), $R(n)$ its operating scope, and $h(n)$ its human accountability anchor. A node $n$ is in autonomous drift if: (1) $\alpha(n) \in \mathcal{N}(t)$ is orphaned from its parent, meaning $\alpha(\alpha(n)) = \text{nil}$ and $\alpha(n) \neq h(n)$; (2) there is no path from $n$ to any human anchor $h$ through iterated parent references; and (3) $n$ continues to execute its event loop and modify state. Drift is not a property of a single node in isolation—it is a property of the chain topology and the operational state of the node simultaneously. It is prevented structurally not by detecting drift at runtime but by making the enabling conditions architecturally impossible.

The Recursive Spawning Protocol (RSP) is the second named pillar of the \bxthree{} Framework, complementing Loop Isolation, Spatial Firewall, Sandbox Gate, and the Bailout Protocol. Together, these five pillars constitute the upstream accountability guarantee: that \bxthree{} systems fail upward into human consciousness, never downward into autonomous chaos. RSP operates across all three \bxthree{} layers in mandatory concert: the Purpose Layer authorizes spawn events and defines the human anchor, the Bounds Engine enforces the depth bound and convergence criteria, and the Fact Layer creates the immutable parent pointer and records every spawn event to an append-only forensic ledger. The protocol makes the delegation chain from every active agent to a human principal a runtime artifact—immutable, cryptographically enforced, and architecturally verifiable—rather than a forensic reconstruction attempted after the fact. There are no advisory phases, no configurable bypasses, and no operational circumstances under which any machine actor may circumvent any protocol requirement.

The remainder of this paper is organized as follows. Section~\ref{sec:rs-background} situates RSP within the research traditions it draws on—agent lifecycle management, accountability and provenance in distributed systems, fixed versus mutable delegation structures, and hierarchical task decomposition—and establishes the \bxthree{} Framework as the minimal sufficient substrate for RSP's accountability guarantees. Section~\ref{sec:rs-drift} provides the formal characterization of autonomous drift, defines the drift triad and its constituent failure sub-modes, catalogs the four concrete failure modes that arise from the three enabling structural conditions, and proves formally that depth limits alone are insufficient to prevent drift. Section~\ref{sec:rs-protocol} specifies the Recursive Spawning Protocol in full: its Purpose Layer validation gate, its Bounds Engine depth management, its Fact Layer immutable pointer and forensic ledger, and its chain verification and convergence criteria. Section~\ref{sec:rs-integration} maps each RSP component to its governing \bxthree{} layer and demonstrates the structural necessity of the three-layer mapping. Section~\ref{sec:rs-correctness} formally states and proves the drift prevention, chain integrity, and termination correctness properties. Section~\ref{sec:rs-limitations} examines the structural costs of RSP's guarantees and identifies directions for future work. Section~\ref{sec:rs-conclusion} concludes with a summary of RSP's contribution to the accountability of multi-agent AI systems.